\section{The lower bound for square radii}\label{sec:square-lower}

\begin{theorem}\label{thm:square-lower}
Let $R_2^*(n)$ be the infimum of the feasible central radii for the
outer radii $1^2,\ldots,n^2$. Then
\[
 \liminf_{n\to\infty}\frac{R_2^*(n)}{n^3}
 \ \ge\ \frac{434+4\sqrt2}{1587\pi}.
\]
\end{theorem}
\begin{proof}
We derive the lower endpoint of \cite[(22)--(23)]{squared} using the
cycles induced by a single feasible placement. If labels $u,v$ have
angular gap $d\in(0,2\pi)$, the cosine law and nonoverlap give
\[
 d\ge2\sin(d/2)\ge
 \frac{2uv}{\sqrt{(R+u^2)(R+v^2)}}\ge \frac{2uv}{R+n^2}.
\]
For a cycle $D$ induced by at least three labels, its consecutive
clockwise gaps sum to $2\pi$. Thus all such cycles in this placement
obey $P(D):=\sum_{\{u,v\}\in E(D)}uv\le\pi(R+n^2)$.

Let $2\le r\le n-2$, put $S=n+r$, and let $D_t$ be the induced cycle
on the labels $t,\ldots,n$. On $D_r$, expanding squares and counting
each vertex twice gives
\[
 P(D_r)=P_{r,n}+\sum_{e\in E(D_r)}\Delta_e,\qquad
 P_{r,n}=\sum_{x=r}^n x(S-x),\quad
 \Delta_{\{u,v\}}=\tfrac12(u+v-S)^2.
\]
Restore labels in decreasing order starting at $r-1$. When $t$ splits
the edge $\{u,v\}$, its signed increment is
$A_t=t(u+v)-uv=t^2-(u-t)(v-t)$.
For $0\le\lambda\le1$ write
\[
 G_\lambda(t)=\frac{\lambda(4St-S^2-2\lambda t^2)}{2(2-\lambda)}.
\]
If the split edge is original, setting $w=u+v$ gives
\[
 \Delta_{\{u,v\}}+\lambda A_t-G_\lambda(t)
 =\frac{\lambda(u-v)^2}{4}
 +\frac{2-\lambda}{4}
 \left(w-\frac{2(S-\lambda t)}{2-\lambda}\right)^2\ge0.
\]
Otherwise one endpoint is a previously restored label at most $r-1$.
Since $u,v>t$, this gives $A_t\ge (S-1)t-n(r-1)$, and
\[
 \lambda((S-1)t-n(r-1))-G_\lambda(t)
 =\frac{\lambda((n-r)^2+4(n-t)+2\lambda(r-1-t)(n-t))}
        {2(2-\lambda)}\ge0 .
\]
Each original edge can be split only once. Its slack therefore pays
for at most one increment, including a split at the cyclic seam.

Choose $r=s_0>s_1>\cdots>s_k\ge1$ and weights
$1\ge\lambda_1\ge\cdots\ge\lambda_k\ge0$. Take the convex combination
of $P(D_{s_0}),\ldots,P(D_{s_k})$ with coefficients
$1-\lambda_1,\lambda_1-\lambda_2,\ldots,\lambda_k$. Each increment in
$s_i\le t<s_{i-1}$ then has coefficient $\lambda_i$.
The difference $G_\lambda(t+1)-G_\lambda(t)
=\lambda(2S-\lambda(2t+1))/(2-\lambda)$ is nonnegative on this range.
Summing the split estimates, discarding unused nonnegative slack, and
taking the infimum over placements yields
\begin{equation}\label{eq:square-finite}
 \pi(R_2^*(n)+n^2)\ \ge\
 P_{r,n}+\sum_{i=1}^k(s_{i-1}-s_i)G_{\lambda_i}(s_i).
\end{equation}

Fix $1/3<\alpha<1/2$, $s=1+\alpha$, $A=3\alpha-1$, and $k\ge1$.
For $0\le i\le k$ set $x_i=1-i/(k+1)$ and
$\beta_i=(s+Ax_i)/4$. Take $\lambda_i=Ax_i/\beta_i$ for $i\ge1$.
Since $s/4<\beta_k<\cdots<\beta_0=\alpha<s/3$ and
$\lambda_i=4-s/\beta_i$, the weights lie in $(0,1)$ in decreasing order.
With $r=\lfloor\alpha n\rfloor$ and $s_i=\lceil\beta_i n\rceil$
for $i\ge1$, the required strict separations hold for all sufficiently
large $n$ at each fixed $k$. Dividing \eqref{eq:square-finite} by $n^3$,
the base term tends to
$p(\alpha)=\int_\alpha^1 x(s-x)\,dx
=(1-\alpha)(\alpha^2+4\alpha+1)/6$. Substitution of
$s=(4-\lambda_i)\beta_i$ shows that
$G_{\lambda_i}(s_i)/n^2\to(\lambda_i\beta_i)^2/2=A^2x_i^2/2$.
As $\beta_{i-1}-\beta_i=A(x_{i-1}-x_i)/4$, we get
\[
 \pi\liminf_n R_2^*(n)/n^3 \ \ge\
 p(\alpha)+\frac{A^3}{8}\sum_{i=1}^k(x_{i-1}-x_i)x_i^2 .
\]
The same liminf satisfies this inequality for each fixed $k$. The sum
is $k(2k+1)/(6(k+1)^2)$ and tends to $1/3$. Choosing
$\alpha=(13-2\sqrt2)/23\in(1/3,1/2)$ gives
\[
 p(\alpha)+A^3/24
 = (23\alpha^3-39\alpha^2+21\alpha+3)/24
 = (434+4\sqrt2)/1587 .
\]
This proves the stated inequality. The power-radius corollary in the
article takes $N_n=n^2$ and marks $(i/n)^2$, so its coefficient $C_2$
is the limit of this same sequence.
\end{proof}
